\SourceLeaf{427}%
the element $\delta\varrho$, then
\[
\delta x=X\delta\varrho,\qquad\delta y=Y\delta\varrho,\qquad\delta z=Z\delta\varrho
\]
and
\[
\delta V=P\delta x+Q\delta y+R\delta z=\mu\delta\varrho.
\]
We thus see how the sign of $\mu$ depends on the change in sign of the value of $V$ when passing from the lower to the upper side.

\Needspace{5\baselineskip}
\begin{center}9.\end{center}
\noindent\hspace*{1.2em}Cut the curved surface by a plane at the point to which our notation refers. This produces a plane curve, of which $ds\mathord{\hbox{\begin{tikzpicture}[x=.1pt,y=-.1pt,baseline=-7pt]
\path[use as bounding box] (0,0) rectangle (17,101);
\path[fill=black] (14,69) -- (12,69) -- (12,70) -- (9,72) -- (9,74) -- (3,80) -- (2,82) -- (2,85) -- (0,87) -- (0,88) -- (2,90) -- (2,96) -- (3,97) -- (3,99) -- (5,100) -- (11,100) -- (14,97) -- (14,88) -- (13,87) -- (13,84) -- (14,83) -- (14,78) -- (16,76) -- (16,75) -- (14,73) -- cycle (6,3) -- (5,4) -- (5,12) -- (6,14) -- (9,14) -- (14,9) -- (14,6) -- (11,3) -- cycle;
\end{tikzpicture}
}}$ is to be an element for which we retain the notation above. We take the side on which the normal to the curved surface lies as the upper side of the plane. We erect a normal to the plane, whose direction is to be denoted by the point $\mathfrak{L}$ on the auxiliary sphere. As we proceed along the curve, $\lambda$ and $L$ therefore change position, while $\mathfrak{L}$ remains constant, with $\lambda L$ and $\lambda\mathfrak{L}=90^{0}$ throughout. Hence $\lambda$ describes the great circle of which $\mathfrak{L}$ is one pole. The element of this great circle will be
\[
=\frac{ds}{r}
\]
if $r$ denotes the radius of curvature of the curve. If the direction of this element is in turn denoted by $\lambda'$ on the auxiliary sphere, then $\lambda'$ clearly lies on the same great circle and is separated from both $\lambda$ and $\mathfrak{L}$ by $90^{0}$. If we now further set
\[
\cos\lambda'(1)=\xi',\qquad\cos\lambda'(2)=\eta',\qquad\cos\lambda'(3)=\zeta',
\]
then
\[
d\xi=\xi'\frac{ds}{r},\qquad d\eta=\eta'\frac{ds}{r},\qquad d\zeta=\zeta'\frac{ds}{r}
\]
since $\xi,\eta,\zeta$ are in fact the coordinates of the point $\lambda$ relative to the center of the sphere.
